% 원제: 15. 가설 충족 현황과 한계
\chapter{가설 충족 현황과 한계}

\section{가설별 현황}

다섯 가설 가운데 네 가설은 해당 모델 범위에서 충족했고, 기준 구조 대비 이점(H5)은 사양 팩에서 기각되어 적용 범위와 산업 방식 대비 성립 조건으로 대체했다(표 \ref{t:status}). 모든 판정은 해당 모델과 조건 안에서만 유효하며, 하드웨어 성능을 뜻하지 않는다.

{\footnotesize
\begin{longtable}{L{\dimexpr 0.1754\linewidth-2\tabcolsep\relax}L{\dimexpr 0.1667\linewidth-2\tabcolsep\relax}L{\dimexpr 0.4298\linewidth-2\tabcolsep\relax}L{\dimexpr 0.2281\linewidth-2\tabcolsep\relax}}
\caption{가설별 현재 상태}\label{t:status}\\
\toprule
\tabhead \textbf{가설} & \textbf{상태} & \textbf{근거} & \textbf{남은 검증} \\
\midrule
\endfirsthead
\toprule
\tabhead \textbf{가설} & \textbf{상태} & \textbf{근거} & \textbf{남은 검증} \\
\midrule
\endhead
\bottomrule
\endlastfoot
H1 정격 분리 & 충족 & 주행 링크 정상 최대 6.44 kW·과도 7.35 kW(평균)·7.19 kW(스위칭), 충전 중 정적 6.68 kW·과도 7.49 kW. 충전전류는 Stage B 비경유(회로 TB07에서 확인) & — \\
H2 강건 제어 & 충족(모델 범위) & 평균모델·독립 구현·스위칭 모델 전부 통과, 회로 대표 사례 일치(TB04–TB08). 창 조정 모듈 공차 포함 456/456·264/264. 실제 Stage A·조정형 Stage C를 넣은 R2 통합 회로 TB16 3/3·TB17 2/2·TB19 2/2 & HIL·시제품 \\
H3 충전 중 보조버스 & 충족(조건 1개) & 동시 운전 Tier-1 460/460·Tier-2 78/78, 모드 전환 240/240, 스위칭 48/48, 회로 한계 7/7(E16, TB07). 충전 중 Stage A 정지 순서와 V\textsubscript{A} 클램프(E26 32/32, TB20 3/4) & 조건: 충전전류 증가 램프 ≥ 100 µs. 영교차 게이트 방향 반전 처리 \\
H4 부하 승인 & 충족(축약 모델) & 예산 초과 0, 승인 지연 최대 17 ms, 미공급 7.0 \%(E8) & 분리·감발 명령 연동, 전압 의존 예산 \\
H5 기준 구조 대비 이점 & 기각(사양 팩) & ③은 ④에 손실·VA·반도체 수 모두 열세(E15) & — \\
Q3 적용 범위·산업 방식 대비 & 조건부 성립 & 창 조정 모듈이 B\textsubscript{420}의 200–216s, B\textsubscript{470}의 224–240s에서 성립 구조 중 최소 손실(E18). 400 V급 보조부품 조건에서 R2가 모터 승압 구조보다 보조 손실 26 \% 작음, 800 V급 부품 설계 대비 이점 없음(E25) & 하드웨어 효율·비용 \\
\end{longtable}
}

\section{현재 근거의 한계}

\begin{itemize}
\item \textbf{모델 수준.} 모든 결과는 시뮬레이션이다. 회로 모델의 스위치 셀은 이상 스위치, 다이오드와 선형 출력용량으로 구성되며, 소자 스위칭 손실, 게이트 드라이버, 배선 기생성분, 열, 절연과 전자기 간섭은 포함하지 않는다.
\item \textbf{손실·정격 비교.} E15·E18·E25는 설계 단계 추정이다(대표 소자, 자성체 비율 규칙, 정적 성립성). 측정된 2:1 하드웨어(전부하 97.4 \%)와 모델(99.7 \%)의 차가 커서 상대 비교로만 쓴다. 모터 승압의 상저항·영상 인덕턴스는 공개 값이 없어 범위로 두었고, 리플 철손·자석 손실을 넣지 않아 모터 승압에 유리하다. 열·기생 성분·비용은 넣지 않았다.
\item \textbf{조정형 Stage C의 하드웨어 값.} 데드타임 창의 폭이 20–30 ns라 GaN의 비선형 출력용량과 게이트 구동 지연 편차로 다시 정해야 한다. 역방향 0.6–0.74 kW에서 이득이 0.96–0.97로 기준을 최대 약 1 \%p 벗어난다. 변압기 권선 용량은 넣지 않았다.
\item \textbf{Stage A.} 통합 회로는 모듈 하나로 계산했다. 모듈 셰딩과 위상 셰딩의 조합, 위상별 공진주파수 공차에서의 셰딩 문턱, 멈춘 위상 플라잉 커패시터의 누설은 확인하지 않았다.
\item \textbf{영교차 게이트의 방향 반전.} 반주기 첨두가 무장 문턱(2.5 A)보다 작은 방향 반전에서 턴오프 전류가 1 A를 조금 넘을 수 있다(TB20 사례 D, 1.01 A). 게이트 개선은 아직 적용하지 않았다.
\item \textbf{충전 중 정지.} 충전기 반응 시간은 CHAdeMO 계열 규격\cite{b16}의 30 ms를 썼다. CCS 계열 규격의 비상 정지 시간은 이 조사에서 값을 확인하지 못했고, 다르면 흡수 에너지가 비례해 바뀐다. 클램프 문턱 여유가 0.6 V 수준이라 비교기·분압기 정확도가 중요하다.
\item \textbf{보호 하드웨어.} V\textsubscript{A} 클램프(7 kJ·220 kW), 링크 제동 저항의 펄스 정격과 반복 고장(고장당 ≤ 22 J), 바이패스 소자의 선형 영역 안전동작영역, 누설 인덕터·변압기 포화 전류(강등 중 링크 쪽 약 70 A 이상), 필터 인덕터 포화 전류(≥ 150 A)는 정격 확인이 남아 있다.
\item \textbf{부하 쪽 응답과 범위.} 400 V 부하가 20 ms 안에 분리·감발한다고 가정했다. 창 조정 모듈의 회로 확인은 R2-W460만 했다(R2-W410은 평균모델과 스위칭 모의). 선행 조사는 12건이며 IEEE Xplore·특허 전수 검색은 하지 않았다.
\end{itemize}
